\chapter*{Conclusion générale}
\addcontentsline{toc}{chapter}{Conclusion générale}

Le projet \projectname{} a permis de concevoir une plateforme de recrutement d'entreprise cohérente, couvrant l'ensemble du cycle métier : centralisation des CV, structuration des offres, progression des candidats avec historique auditable, organisation des entretiens, notifications, analytique, gouvernance des actions IA et administration. Contrairement à une application limitée au suivi administratif, la solution rassemble dans un même environnement les opérations, les données et les mécanismes de décision qui structurent réellement le travail des équipes TA, Manager, HR et Admin.

Sur le plan technique, le projet montre la pertinence d'une architecture \textit{feature-driven} pour un produit full-stack moderne. La séparation entre routes, actions serveur, services métier, composants d'interface et schéma de données favorise la lisibilité, la maintenabilité et l'évolution incrémentale. Cette organisation a également facilité l'intégration de briques avancées comme les embeddings, les chunks RAG, les guides d'entretien IA et l'agent conversationnel multi-outils, sans dénaturer le c\oe{}ur métier de l'application.

La contribution la plus distinctive réside dans l'industrialisation raisonnée de l'intelligence artificielle. Le projet ne s'appuie pas sur un modèle unique utilisé de manière opaque ; il combine extraction structurée, matching, génération guidée, autopilot d'entretien, conversation outillée, garde-fous d'exécution et validation des sorties. L'évaluation comparative du pipeline RAG sur 40 requêtes, avec 85\% de couverture de vérité terrain, montre un gain de pertinence mesurable : amélioration de Precision@5, Precision@10, MRR@10 et NDCG@10, tout en maintenant un taux d'erreur nul.

\section*{Validation fonctionnelle et technique}
La validation retenue pour ce rapport combine deux niveaux complémentaires : d'une part une validation fonctionnelle par scénarios couvrant les principaux parcours utilisateurs du produit, et d'autre part une validation technique fondée sur les preuves automatisées documentées dans les artefacts du projet. Cette approche permet de ne pas limiter la démonstration à la seule couche IA, mais d'élargir la vérification au fonctionnement global de la plateforme. Le tableau \ref{tab:functional-validation} synthétise cette validation fonctionnelle.

\begin{table}[H]
  \centering
  \scriptsize
  \caption{\label{tab:functional-validation}Synthèse de validation fonctionnelle}
  \setlength{\tabcolsep}{3pt}
  \noindent\begin{adjustbox}{max width=\linewidth,center}
  \begin{tabularx}{\textwidth}{>{\bfseries\raggedright\arraybackslash}p{2.4cm}>{\raggedright\arraybackslash}X>{\raggedright\arraybackslash}X>{\raggedright\arraybackslash}X}
    \toprule
    Domaine & Scénario de validation & Résultat attendu & Résultat retenu \\
    \midrule
    Accès et identité & Connexion puis accès à l'espace métier. & Redirection par rôle et blocage sans session valide. & Couvert par l'authentification, la redirection par rôle et le shell dashboard. \\
    Catalogue de CV & Import d'un CV PDF ou DOCX. & Extraction du texte, profil consultable et enrichissement lancé. & Couvert par l'upload, l'extraction, les doublons et l'indexation progressive. \\
    Pipeline candidat & CV affecté puis transmis entre TA, Manager et RH. & Candidat créé, transitions historisées et responsabilités tracées. & Couvert par le workflow candidat, les affectations et l'historique d'étapes. \\
    Entretiens et décisions & Entretien planifié puis formalisé par un rapport. & Entretien créé, étape mise à jour et onboarding ouvert si embauche. & Couvert par la planification, le reporting et la checklist d'onboarding. \\
    Administration et gouvernance & Consultation des vues admin et du centre de gouvernance. & Supervision consolidée, traces d'audit, confirmations IA et exports. & Couvert par les vues admin, les logs, les exports et le centre de gouvernance. \\
    Chat analytique & Question posée depuis le chat ou \sourcecode{/agent}. & Outils RBAC, streaming SSE, preuves, confirmations, graphiques, cartes typées et trace par phases. & Couvert par l'orchestration, le grounding, les confirmations, les cartes de réponse, les graphiques et la timeline d'outils. \\
    \bottomrule
  \end{tabularx}
  \end{adjustbox}
\end{table}

\paragraph{Validation technique automatisée.}
Sur le plan technique, les contrôles automatisés exécutés pour cette révision ont validé \sourcecode{bun run lint} sans avertissement et \sourcecode{bun run test} : 27 fichiers de test et 125 tests Vitest passent. La compilation de production \sourcecode{bun run build} fait partie de la vérification finale avant livraison.

Les vérifications automatisées couvrent notamment :
\begin{itemize}
  \item le grounding des noms candidats, avec rejet des noms absents des résultats d'outils, acceptation des noms autorisés et remplacement déterministe des classements hallucinés ;
  \item l'isolation de périmètre, avec filtrage des résultats pour les utilisateurs non administrateurs, prévention des fuites inter-utilisateurs et séparation des clés de cache selon l'utilisateur et le rôle ;
  \item le comportement du chat en cas d'absence de résultats exploitables, avec réponse de repli grounded au lieu d'une fabrication libre ;
  \item la confirmation des actions IA modifiant les données, depuis la création d'une action en attente jusqu'à son exécution, son annulation ou son expiration ;
  \item les transitions du pipeline candidat, avec validation des transitions autorisées, rejet des transitions incohérentes et formatage de la timeline ;
  \item la génération des métadonnées d'évidence, des liens profonds et des résumés de gouvernance administrateur ;
  \item la qualification des panneaux d'\textit{evidence readiness} côté candidat et côté catalogue documentaire ;
  \item la génération, la sérialisation et l'extraction des graphiques du chat analytique, afin que les réponses visuelles restent fondées sur des sorties d'outils validées ;
  \item la génération, la sérialisation et l'extraction des cartes de réponse typées du chat, afin que les réponses candidat, pipeline et gouvernance soient rendues à partir de sorties d'outils validées ;
  \item le regroupement de la trace d'outils par phases, avec classification en récupération, analyse, confirmation, exécution et vérification ;
  \item le centre de gouvernance administrateur, incluant l'affichage des lignes d'audit et l'ouverture de détails sanitizés.
\end{itemize}

Enfin, la couche de retrieval fait l'objet d'une validation quantitative spécifique déjà documentée dans le projet. Le rapport d'évaluation RAG de phase 2 couvre 40 requêtes avec 85\% de couverture \textit{ground truth} et montre une amélioration mesurable de Precision@5, Precision@10, MRR@10 et NDCG@10. Cette mesure complète utilement la validation fonctionnelle, car elle montre non seulement que le produit fonctionne, mais aussi que sa couche de recherche augmentée est évaluée de manière chiffrée.

\section*{Atteinte des objectifs initiaux}
L'évaluation finale doit être reliée aux objectifs formulés au départ du projet. Le tableau \ref{tab:objective-achievement} synthétise cette lecture : il ne s'agit pas seulement de vérifier que des écrans existent, mais de montrer que chaque objectif produit une preuve fonctionnelle, technique ou expérimentale.

\begin{table}[H]
  \centering
  \scriptsize
  \caption{\label{tab:objective-achievement}Évaluation de l'atteinte des objectifs initiaux}
  \begin{adjustbox}{max width=\linewidth,center}
  \begin{tabularx}{\textwidth}{>{\bfseries\raggedright\arraybackslash}p{2.8cm}>{\raggedright\arraybackslash}X>{\raggedright\arraybackslash}X>{\raggedright\arraybackslash}p{2cm}}
    \toprule
    Objectif initial & Preuve de réalisation & Limite identifiée & Statut \\
    \midrule
    Centraliser les données RH & CV, offres, candidats, entretiens, notifications, logs et onboarding sont regroupés dans une base PostgreSQL commune. & La qualité dépend encore de la complétude des CV et des champs renseignés. & Atteint \\
    Structurer le pipeline candidat & États explicites, affectations TA/Manager/HR, rapports et historique immuable des transitions. & Des scénarios end-to-end production seraient utiles pour mesurer les délais réels par étape. & Atteint \\
    Accélérer le matching & Matching explicite, recherche sémantique, hybrid search et bulk assignment. & Le RAG avancé améliore la pertinence mais augmente fortement la latence. & Partiellement optimisé \\
    Rendre l'IA explicable & Sources, grounding, cartes typées, trace outillée, confirmations et résultats RAG mesurés. & L'explicabilité doit être testée avec des utilisateurs métier réels. & Atteint \\
    Gouverner les actions sensibles & RBAC, validation Zod, pending actions, logs d'activité, gouvernance admin et exports. & Une politique de rétention et d'audit long terme reste à définir selon les contraintes organisationnelles. & Atteint \\
    Produire une base évolutive & Architecture feature-driven, services métier, composants réutilisables et tests automatisés. & Les modules IA devront être profilés avant passage à un volume de données plus important. & Atteint avec réserves \\
    \bottomrule
  \end{tabularx}
  \end{adjustbox}
\end{table}

Cette lecture montre que les objectifs fonctionnels majeurs sont atteints, tandis que l'optimisation de performance et la validation métier en conditions réelles restent les deux axes les plus importants pour transformer le prototype avancé en solution industrialisée.
\section*{Discussion des résultats}
Les résultats observés permettent de tirer plusieurs enseignements utiles au-delà du simple constat chiffré. D'abord, l'amélioration de la pertinence du pipeline RAG s'explique par le passage d'une récupération trop limitée à une récupération hybride plus riche. La combinaison entre réécriture de requête, recherche vectorielle, recherche lexicale, fusion par RRF, reranking et assemblage de contexte donne au système davantage de chances de retrouver des passages réellement alignés avec l'intention métier de la requête.

La chaîne de retrieval enrichie donne au système une lecture plus complète de la demande utilisateur. Chaque étape apporte un signal distinct : la reformulation clarifie l'intention, la récupération vectorielle capte les proximités sémantiques, la recherche lexicale préserve les termes exacts, la fusion réduit la dépendance à un seul score et le reranking réordonne les résultats selon leur utilité finale.

Les mécanismes de grounding jouent ici un rôle déterminant. Sans eux, le gain de fluidité conversationnelle pourrait se payer par des réponses plus séduisantes que fiables. Le grounding garantit que les noms candidats, les classements et les restitutions restent bornés par les résultats réellement accessibles à l'utilisateur. Dans un contexte RH, cette contrainte est essentielle, car une réponse plausible mais non fondée peut dégrader la confiance dans l'outil et produire des erreurs de décision.

Pour le recruteur, l'impact réel de ces résultats est double. D'un côté, la plateforme réduit l'effort de recherche, de comparaison et de préparation grâce à une récupération plus pertinente et à un accès conversationnel outillé. De l'autre, elle conserve une logique de contrôle compatible avec un usage professionnel : les réponses sont contextualisées, les outils sont filtrés par rôle et les sorties sensibles sont vérifiées. Le résultat n'est donc pas seulement une IA plus performante, mais une IA plus exploitable dans un processus de recrutement gouverné.
L'ajout de visualisations directement dans les réponses agentiques renforce ce résultat côté usage. Les questions analytiques ne produisent plus uniquement une synthèse textuelle : lorsque les données s'y prêtent, elles affichent une courbe ou des barres qui rendent les tendances, les goulets d'étranglement et les écarts de compétences plus immédiatement compréhensibles.
La confirmation explicite des actions modifiant les données complète cette logique. Elle évite qu'une réponse agentique transforme directement l'état du système sans validation humaine, tout en conservant une trace exploitable par le centre de gouvernance.
Les cartes de réponse typées et la trace par phases prolongent la même logique d'explicabilité. Elles réduisent l'effort de lecture : les métriques, liens, limites et actions de suivi apparaissent sous forme structurée, tandis que le travail interne de l'agent est présenté comme une timeline métier plutôt que comme une liste brute d'appels techniques.

\section*{Vue d'ensemble des acquis}
La vue d'ensemble du projet met en évidence une progression cohérente entre les objectifs initiaux, le Product Backlog et les sprints de réalisation. Le produit démarre par un accès sécurisé, construit ensuite ses objets métier principaux, organise le workflow candidat, ajoute une couche IA gouvernée puis termine par des capacités de supervision. Cette continuité donne au rapport une lecture claire : chaque module apporte une valeur immédiate tout en préparant les modules suivants.

Le résultat obtenu est une plateforme unifiée où les données RH ne sont plus dispersées entre des écrans isolés. Les CV, les offres, les candidats, les entretiens, les notifications, les logs et les analyses IA s'inscrivent dans un même cycle métier. Cette vue globale est essentielle, car elle montre que \projectname{} n'est pas seulement une démonstration technique ; c'est un système orienté usage, pensé pour aider les équipes à suivre, comprendre et justifier leurs décisions.

\section*{Limites du projet}
Malgré les résultats obtenus, plusieurs limites doivent être soulignées. La première concerne la latence du pipeline RAG avancé : l'amélioration de la pertinence s'accompagne d'un coût de calcul important lorsque la requête déclenche l'ensemble de la chaîne de retrieval enrichie. La deuxième limite tient à la dépendance envers des services externes de génération et d'embeddings, qui peuvent introduire des contraintes de disponibilité, de coût et de variabilité de performance.

Une autre limite concerne la qualité des données d'entrée. La valeur produite par le matching, le screening, le RAG et l'agent conversationnel dépend directement de la qualité des CV importés, de la cohérence des offres d'emploi et de la richesse des informations renseignées dans le pipeline candidat. Enfin, la validation automatisée observable couvre désormais le grounding, l'isolation de périmètre, les confirmations agentiques, les transitions candidat, les panneaux d'évidence, les cartes de réponse typées, la trace d'outils par phases, le centre de gouvernance et l'évaluation RAG ; elle gagnerait encore à être complétée par davantage de scénarios end-to-end métier exécutés dans un environnement proche production.

\section*{Menaces à la validité et risques éthiques}
Dans un projet de recrutement augmenté par l'IA, la validation ne peut pas se limiter à la réussite technique des fonctionnalités. Les résultats doivent être lus avec prudence, car ils dépendent du corpus de CV, des requêtes choisies, du comportement des utilisateurs et du cadre éthique dans lequel l'outil est utilisé. Cette précaution est particulièrement importante dans le domaine RH, où une recommandation algorithmique peut influencer une décision humaine sensible. Les travaux sur les algorithmes de recrutement soulignent d'ailleurs que les promesses de réduction du biais doivent être évaluées de manière critique et non acceptées comme un effet automatique de l'automatisation \netref{web:raghavan-hiring}.

\begin{table}[H]
  \centering
  \scriptsize
  \caption{\label{tab:validity-ethics-risks}Menaces à la validité et mesures de maîtrise éthique}
  \begin{adjustbox}{max width=\linewidth,center}
  \begin{tabularx}{\textwidth}{>{\bfseries\raggedright\arraybackslash}p{2.6cm}>{\raggedright\arraybackslash}X>{\raggedright\arraybackslash}X>{\raggedright\arraybackslash}X}
    \toprule
    Dimension & Menace ou risque & Impact possible & Mesure de maîtrise \\
    \midrule
    Validité interne & Les résultats RAG dépendent du jeu de données local, des chunks générés et des requêtes retenues. & Une amélioration mesurée peut refléter le corpus testé plus qu'une supériorité générale. & Comparaison baseline/phase 2 sur le même protocole, conservation des métriques et extension future du jeu de requêtes. \\
    Validité externe & Le système n'est pas encore évalué dans un environnement de production avec plusieurs recruteurs réels. & Les gains observés peuvent différer selon les pratiques métier, les volumes et les profils traités. & Prévoir une expérimentation utilisateur avec TA, Manager et HR, puis mesurer temps, satisfaction et qualité perçue. \\
    Validité de construction & Precision@K, MRR et NDCG mesurent la pertinence de retrieval, pas directement la qualité d'une décision RH. & Un bon score informationnel ne garantit pas une décision juste ou acceptable métier. & Compléter les métriques IR par feedback utilisateur, audit des décisions et indicateurs de qualité d'entretien. \\
    Validité statistique & Le protocole couvre 40 requêtes avec 85\% de vérité terrain. & La variance reste possible, surtout par catégorie de requête ou niveau de difficulté. & Augmenter le nombre de requêtes, équilibrer les catégories et documenter les intervalles de confiance. \\
    Biais algorithmique & Les CV et historiques peuvent contenir des biais sociaux, linguistiques ou de parcours. & Le score IA pourrait reproduire ou amplifier des déséquilibres existants. & Conserver une décision humaine, afficher les critères, auditer les écarts et éviter les décisions automatiques finales. \\
    Hallucination et opacité & Un agent génératif peut citer un candidat absent ou formuler une justification non fondée. & Perte de confiance et risque de décision sur information incorrecte. & Grounding des candidats, citations, sources, cartes typées, trace d'outils et fallback lorsque les données sont insuffisantes. \\
    Données personnelles & Les CV contiennent des informations sensibles ou indirectement identifiantes. & Risque de fuite, de consultation hors rôle ou de conservation excessive. & RBAC, séparation des rôles, journaux d'activité, minimisation des sorties et politique future de rétention. \\
    Sur-confiance utilisateur & Un score ou une synthèse IA peut être perçu comme plus objectif qu'il ne l'est réellement. & L'utilisateur peut déléguer trop vite son jugement au système. & Présenter les limites, les gaps, les preuves et les recommandations comme aide à la décision, jamais comme décision automatique. \\
    \bottomrule
  \end{tabularx}
  \end{adjustbox}
\end{table}

Le tableau \ref{tab:validity-ethics-risks} montre que la contribution du projet ne réside pas seulement dans l'ajout de fonctions IA, mais aussi dans leur encadrement. Les mécanismes de grounding, de confirmation, de RBAC, de logs et d'explication ne sont donc pas secondaires : ils conditionnent la possibilité d'utiliser l'IA dans un contexte RH sans transformer l'assistance en automatisation opaque.

\section*{Perspectives}
Plusieurs prolongements peuvent être envisagés. La première perspective consiste à améliorer les métriques et la mesure de complexité : augmenter le nombre de requêtes d'évaluation, distinguer les métriques par type de besoin métier, suivre la latence par étape du pipeline, mesurer le coût financier par requête et intégrer des indicateurs de satisfaction utilisateur. Cette évolution permettrait de comparer objectivement plusieurs stratégies : SQL seul, vectoriel seul, RAG hybride, RAG avec reranking et agent multi-outils.

La deuxième perspective concerne l'optimisation du pipeline RAG. Les pistes prioritaires sont la réduction du \textit{topK} lorsque la requête est simple, le cache des embeddings et des résultats, l'activation conditionnelle de la réécriture, la mesure séparée du coût de reranking et l'étude d'index vectoriels plus adaptés aux volumes importants. Un fine-tuning ciblé peut aussi être envisagé, non pour mémoriser les CV, mais pour stabiliser des tâches répétitives : classification d'intention, extraction d'exigences, reformulation de requêtes ou normalisation des compétences.

La troisième perspective porte sur le recrutement prédictif. À partir d'un historique suffisant et gouverné, la plateforme pourrait estimer les risques de blocage dans le pipeline, prédire les délais de passage entre étapes, recommander les profils avec une probabilité de réussite contextualisée et détecter les écarts entre demande métier et disponibilité du catalogue. Cette évolution devra cependant rester explicable et contrôlée, car une prédiction RH non justifiée peut introduire des biais de décision.

Enfin, l'intégration de modèles multimodaux représente une perspective de recherche intéressante. Elle permettrait d'exploiter non seulement le texte des CV et offres, mais aussi des documents scannés, portfolios, certificats, tableaux ou supports d'entretien. Cette ouverture doit être accompagnée de règles de gouvernance renforcées : consentement, minimisation des données, filtrage des informations sensibles et séparation claire entre aide à l'analyse et décision humaine.
\section*{Apports personnels}
Au regard des réalisations présentées dans ce mémoire, les apports personnels du travail mené peuvent être synthétisés autour de plusieurs axes complémentaires. Le premier apport réside dans la structuration d'une architecture cohérente, capable d'articuler authentification, services métier, persistance, interfaces multi-rôles et composants IA dans un même produit full-stack.

Le deuxième apport concerne l'intégration d'une couche d'intelligence artificielle non cosmétique, directement reliée aux données du système : extraction de CV, matching, pipeline RAG hybride, génération de guides, screening et chat analytique outillé. Le troisième apport porte sur la formalisation du workflow de recrutement lui-même, avec une progression explicite des candidats, une répartition claire des responsabilités entre TA, Manager et RH, un historique immuable des transitions, ainsi qu'une continuité jusqu'à l'onboarding et à la gouvernance. Enfin, ce travail a permis de relier l'implémentation technique à une logique de contrôle, grâce à la validation des entrées, au filtrage RBAC, au grounding des réponses, à la confirmation des actions IA sensibles et à la traçabilité des actions critiques.

La dernière amélioration d'interface renforce cet apport en transformant le chat agentique en surface de décision plus lisible : les cartes typées structurent les résultats, la timeline par phases rend les appels d'outils auditables et la suite de tests agentiques stabilise ces comportements.


Cette progression confirme que le travail réalisé couvre les dimensions principales attendues d'une plateforme de recrutement augmentée : sécurité d'accès, centralisation des données, workflow opérationnel, assistance IA, validation quantitative et gouvernance. L'ensemble forme un produit cohérent, documenté et aligné avec les besoins des rôles TA, Manager, HR et Admin.

En définitive, \projectname{} constitue une base solide pour un système de recrutement intelligent orienté entreprise. Le projet démontre qu'il est possible d'intégrer l'IA de manière utile, gouvernée et mesurable, à condition de l'ancrer dans les données du métier, dans des workflows explicites et dans une architecture technique disciplinée.
